\section{Methods}
\label{sec:methods}

\subsection{Postprocessing methods compared}
\label{sec:methods-list}

We compare six families of methods, all sharing a single interface: given a
station's ensemble forecast summary for a lead time, produce a predictive
distribution for the observed \(2\,\mathrm{m}\) temperature at that lead,
expressed as the \NQuantileLevels{} quantiles \QuantileLevelsList{} (a fixed
global constant used identically by every method in this study).

\textbf{Raw ensemble.} The empirical distribution of the raw NWP ensemble
members at each station and lead time, used unmodified. This is the
zero-shot, zero-parameter baseline every postprocessing method is expected to
improve on.

\textbf{EMOS.} Ensemble Model Output Statistics \citep{gneiting_2005_emos}, a
nonhomogeneous Gaussian regression fit by minimizing the closed-form Gaussian
CRPS \citep{gneiting_2007_proper_scoring} over the training set:
\(\mu = a + b \cdot \bar{x}\), \(\sigma^2 = \exp(c) + \exp(d) \cdot s^2\),
where \(\bar{x}\) and \(s^2\) are the ensemble mean and variance. We fit two
variants of this same model, differing only in how many stations share one
parameter set: \emph{pooled} EMOS fits one global \((a,b,c,d)\) across all
\NStations{} stations, following the field's standard practice
\citep{gneiting_2005_emos,raftery_2005_bma}; \emph{local} EMOS fits a
separate \((a,b,c,d)\) per station whenever that station has enough training
rows to do so (we require at least 5 rows; stations below this threshold are
excluded from local EMOS's evaluation at that training size, and the
resulting station-coverage fraction is reported alongside every local-EMOS
result). \citet{baran_2016_combining} and \citet{gebetsberger_2018_estimation}
report comparable station-adaptive EMOS variants at small training sizes,
the regime this paper's breakpoint curve targets directly.

\textbf{DRN.} The Distributional Regression Network of
\citet{rasp_2018_drn}: a feed-forward neural network with a learned
per-station embedding, taking the ensemble mean and variance as input and
outputting \((\mu, \log\sigma^2)\) parameters of the same Gaussian family
EMOS uses, trained by gradient descent (Adam) on the same closed-form
Gaussian CRPS loss, reimplemented in PyTorch so it is differentiable. Unlike
the original DRN, which additionally uses auxiliary NWP fields and lead time
as inputs, our training pipeline pools all \NLeadTimes{} lead times into one
fit and does not expose lead time to any model (Section~\ref{sec:limitations});
DRN and EMOS are therefore compared under the same feature constraint, not
DRN's originally-intended full feature set. \citet{schulz_2021_windgusts}
runs a systematic comparison of neural postprocessing architectures under
this same feature-constrained setup, the closest methodological precedent
to our own DRN evaluation. Time-series/lead-time-continuous EMOS variants
\citep{jobst_2024_ts_emos,wessel_2024_leadtime_continuous} target the
small-training-data and lead-time-resolved regimes this paper's breakpoint
curve also explores, without themselves being baselines here (a wider
survey of the neural-postprocessing and quantile-regression literature is
given in Section~\ref{sec:intro}).

\textbf{TimesFM-3 (zero-shot).} The 3.0 release of TimesFM, a frozen,
pretrained decoder-only time series foundation model. The architecture
family is described in \citet{das_2024_timesfm_icml}; that paper describes the
original TimesFM 1.0 decoder-only architecture and does not cover the
3.0 release's covariate support, which Section~\ref{sec:covariates}
sources directly rather than through that citation. This project uses
checkpoint \texttt{google/timesfm-3.0-pytorch}, with no fine-tuning,
task-specific training, or parameter updates of any kind.
Section~\ref{sec:covariates} describes its input construction. Contrasted
with generalist AI weather-prediction models that operate on full spatial
fields \citep{lam_2023_graphcast,price_2024_gencast,bodnar_2025_aurora} --
TimesFM-3 is used here in univariate, single-station mode, not as a
spatial forecasting system, and is not compared against those models
here.

\textbf{Variance-inflation baseline.} A trivial, zero/low-parameter
alternative: scale the raw ensemble's per-instance spread by a single
multiplier \(\lambda\) and read off Gaussian quantiles at the same
\NQuantileLevels{} levels. We report \NVarianceInflationVariants{} variants,
each answering a different question about whether TimesFM-3's apparent skill
is trivially reproducible \citep{zhang_2025_context_parroting}: (a) a
CRPS-optimal \(\lambda\) fit by minimizing training-set CRPS at each training
size \(N\) (fair comparison against EMOS, which also uses \(N\) training
cases); (b) a fixed climatological \(\lambda\), requiring no data at all
(fair comparison against TimesFM-3's own zero-shot status); (c) a
coverage-matched \(\lambda\), fit on training data only (never on test-set
coverage) by root-finding to reach as closely as a training-fit multiplier
allows TimesFM-3's own empirical coverage, then evaluated for interval
width on the test set -- this directly asks whether TimesFM-3's
calibration, not just its CRPS, is achievable by a one-parameter
rescaling. Because \(\lambda\) is fit on training data and applied to a
different test set, the coverage this variant actually reaches on test
does not land exactly on the training-time target; we report both values
together (Section~\ref{sec:results-full}) rather than the word
``exactly.'' (d) a leakage-optimistic \(\lambda\), fit by root-finding
directly on \emph{test-set} coverage (never on training data) to reach
TimesFM-3's own empirical coverage -- the deliberately
leakage-contaminated counterpart to (c), answering not "what can a
legitimate, leakage-free trivial rescaling achieve" but "what is the
ceiling on what any trivial rescaling could ever achieve at this
coverage target, even with an unfair advantage." Variant (d) is never
presented as a legitimate baseline result; it exists solely to bound how
much of variant (c)'s width disadvantage against TimesFM-3
(Section~\ref{sec:results-full}) could in principle be a leakage
artifact of the train/test split rather than a real limit of trivial
rescaling.

\subsection{TimesFM-3 covariate construction, and the information asymmetry it creates}
\label{sec:covariates}

TimesFM-3's target series is a station's own observation history; its
past-future (dynamic) covariate channel\footnote{Native past-future
covariate support is documented for the 3.0 release specifically, not in
the original TimesFM architecture paper \citep{das_2024_timesfm_icml}: model
card at \url{https://huggingface.co/google/timesfm-3.0-pytorch} and
release notes at
\url{https://research.google/blog/timesfm-3-a-zero-shot-foundation-model-for-multivariate-forecasting/}
(Google Research; both accessed 2026-09-07).} is filled with the ensemble mean,
supplied as a value known across the entire forecast horizon rather than
only in the past. This study supplied only the ensemble mean to this
channel; ensemble spread was not. This is an input-design choice, not a
confirmed architectural limit: the installed package's own source
documents the channel's accepted shape as genuinely multi-channel, and a
preliminary test confirms the model's output is measurably sensitive to
a second channel's real informational content, not merely its presence
(Section~\ref{sec:limitations} reports this preliminary test;
Section~\ref{sec:discussion-covariate-check} reports the full evaluation
it motivated). The model is entirely frozen: no fine-tuning,
no gradient updates, and no task-specific training of any kind are
performed at any point in this study.
\citet{kreusel_2026_covariates_key} independently reports that covariate
handling is the dominant factor in time series foundation model
performance; a coding error encountered during this project's own
development, where an incoherent covariate channel measurably degraded
TimesFM-3's apparent skill before being corrected, is consistent with that
finding (Section~\ref{sec:pitfall}).

This covariate construction leaves TimesFM-3 and the trained statistical
methods (EMOS, DRN) with genuinely different information, as this study
was actually run: TimesFM-3 sees the station's own recent observation
history but was not given ensemble spread; EMOS and DRN see ensemble
spread but no observation history at all (Table~\ref{tab:info-channels}).
We treat this asymmetry as a candidate explanation for three separate
empirical results in this paper (Section~\ref{sec:results}): TimesFM-3's
short-lead advantage (recent observations carry strong persistence
information at short lead times, information EMOS/DRN were not given
observation history to use in this study's configuration),
its calibration deficit relative to EMOS (TimesFM-3's covariate channel
was not given ensemble spread, the direct measure of forecast
uncertainty, in this study), and its long-lead disadvantage (observation
history loses predictive value as lead time grows, leaving TimesFM-3's
remaining information -- the ensemble mean alone -- no richer than what
EMOS already uses). This explanation was conditional on whether ensemble spread could be
added to TimesFM-3's covariate channel without degrading it -- a
question this project has since tested directly
(Section~\ref{sec:discussion-covariate-check}). The channel does accept
a second, simultaneous input (Section~\ref{sec:limitations}), but
supplying spread this way did not close the calibration deficit: CRPS
and interval width both worsen, and coverage moves only marginally
closer to the nominal target at a substantially wider interval. This
does not confirm the information-asymmetry framing for this half of the
mechanism question, but it does not confirm an architectural explanation
either (Section~\ref{sec:discussion-covariate-check},
Section~\ref{sec:conclusion}).
Section~\ref{sec:discussion} tests both halves of this explanation
directly: the observation-history half by augmenting EMOS with a
persistence predictor, and the calibration half by supplying TimesFM-3
with ensemble spread as a second covariate.

\begin{table}[htbp]
\centering
\caption{Information supplied to each method in this study -- not each
method's inherent architectural capability. TimesFM-3's ``no'' for
ensemble spread reflects this study's covariate design, not a confirmed
limit of the 3.0 release (see note below and Section~\ref{sec:limitations}).}
\label{tab:info-channels}
\begin{tabular}{lccc}
\toprule
Method & Observation history & Ensemble mean & Ensemble spread \\
\midrule
EMOS (pooled / local) & no & yes & yes \\
DRN & no & yes & yes \\
TimesFM-3 (zero-shot) & yes & yes & no$^{\ast}$ \\
Raw ensemble & no & yes (empirical) & yes (empirical) \\
Variance-inflation baseline & no & no & yes (scaled) \\
\bottomrule
\end{tabular}

\smallskip
{\small $^{\ast}$Not supplied in this study. TimesFM-3's covariate
parameter accepts a second, simultaneous channel by design (see the note
above), and a full evaluation supplying it found CRPS and interval width
both worsen, not improve (Section~\ref{sec:discussion-covariate-check}).}
\end{table}

Two other covariate-aware zero-shot approaches make a different design
choice than the frozen, built-in covariate channel this paper relies on:
\citet{auer_2025_icl_covariates} inject covariates through in-context
examples rather than an architectural slot, requiring no model changes
but also no guarantee the model attends to them the way a dedicated
covariate channel would; \citet{qin_2025_cora} instead adapt a pretrained
TSFM's covariate handling per-task through a small tuning step, which
this paper's zero-shot, no-fine-tuning design (Section~\ref{sec:methods-list})
deliberately avoids.

\subsection{Metrics}
\label{sec:metrics}

The Continuous Ranked Probability Score (CRPS) is our primary metric.
Every reported CRPS value in this paper, for every method, is computed by
the same pinball-loss-based quantile estimator
\citep{jordan_2017_scoringrules} applied to the identical
\NQuantileLevels{}-level quantile representation: EMOS and DRN's own
closed-form Gaussian CRPS \citep{gneiting_2007_proper_scoring,hersbach_2000_crps}
is used only inside their training objective (Section~\ref{sec:methods-list}),
never for a reported evaluation metric, so no method receives an
analytic-CRPS advantage over the discretized estimator every other method
is necessarily scored with. We additionally report the tail-weighted CRPS
(twCRPS), the same pinball-loss estimator restricted to the upper two
quantile levels (0.8, 0.9) \citep{allen_2023_weighted_scoringrules,taillardat_2019_extremes},
as a check on tail behavior; \citet{zamo_2018_crps_estimation},
\citet{bolin_2019_scale_invariance}, and \citet{arnold_2023_crps_decomposition}
discuss estimator choice and decomposition properties of CRPS more broadly,
and \citet{leutbecher_2020_crps_changes} interprets CRPS differences
directly. We additionally report \emph{coverage@80\%}: the empirical
fraction of observations falling within the predicted p10--p90 interval,
against a nominal level of 0.80 -- not 0.90, since our fixed
\NQuantileLevels{}-level quantile grid has no p05/p95 pair to compute a
90\%-nominal interval from; every mention of coverage in this paper refers
to this 80\%-nominal band and is labeled accordingly throughout. Finally we
report \emph{interval width} (the p10--p90 range in Kelvin) as a measure of
sharpness. Per \citet{gneiting_2007_calibration_sharpness}'s principle of
maximizing sharpness subject to calibration, interval width is reported as a
curve for every method and is never assigned a single ``better'' direction
or a breakpoint of its own (Section~\ref{sec:breakpoint-def}): a narrower
interval is only meaningful jointly with adequate coverage, not on its own.
PIT histograms (Section~\ref{sec:results}) are likewise reported
descriptively; \citet{broecker_2018_serial_dependence} and
\citet{broecker_2020_stratified_rank} show that formal rank-uniformity
hypothesis tests require correction for serial dependence, which our PIT
values genuinely exhibit, so no such test is applied here
(Section~\ref{sec:limitations}). \citet{dirkson_2025_misdiagnosing_reliability}
further shows that coverage and rank-based reliability diagnostics can both
be misled by the same covariance pathology; we do not claim these
descriptive diagnostics are immune to that critique (Section~\ref{sec:limitations}).

\subsection{Statistical significance}
\label{sec:significance}

The evaluation set is a station~$\times$~time~$\times$~lead panel, not a
single independent sample: instances share temporal autocorrelation within
a station and spatial (synoptic) dependence across the \NStations{}
stations on a given day. A textbook Diebold--Mariano test
\citep{diebold_1995_comparing_accuracy,diebold_2015_twenty_years} assumes one
autocorrelated series and is not directly applicable to this panel structure;
its small-sample behavior is separately unreliable even in the single-series
case it was designed for \citep{coroneo_2020_small_samples}. We test two
block definitions, aggregating the per-instance loss differential to one
mean value per block before testing: a \emph{station block} (absorbing
within-station temporal correlation, but not cross-station dependence) and
a \emph{day block} (absorbing cross-station synoptic dependence on a given
day, following the logic of exact tests under contemporaneous correlation,
\citealp{luger_2004_exact_tests}, but not multi-day persistence in the loss
differential itself). We do not additionally test a multi-day moving-block
variant that would absorb such short-range temporal persistence; the day
block already targets the specific mechanism this analysis was designed to
address (same-day synoptic systems affecting many stations at once), and we
report this as a remaining, open block-definition question rather than as a
claim that day-blocking is fully sufficient. For each block definition we
report a paired \(t\)-test and a Wilcoxon signed-rank test on the block
means, and a block-bootstrap confidence interval on the CRPS skill score
(resampling whole blocks with replacement). We call this a
\emph{station-blocked} or \emph{day-blocked paired test}, not a
Diebold--Mariano test: the classical test's assumptions do not transfer to
this station~$\times$~time~$\times$~lead panel.

\subsection{Breakpoint definition}
\label{sec:breakpoint-def}

The paper's central quantity is the training-data size at which a trained
method's performance crosses a zero-shot reference's. We define \(N\) days
of training as \(k = \mathrm{round}(N / \DaysPerCase)\) issue-date cases,
\DaysPerCase{} being the reforecast archive's measured ratio of calendar
days to sampled cases (Section~\ref{sec:data}). A breakpoint is
the linearly interpolated case count \(k\) at which a trained method's
metric crosses the zero-shot reference's value, computed over an ordered
sequence of tested \(k\); we distinguish three outcomes explicitly rather
than collapsing them into one ``no breakpoint'' label, which conflates
opposite situations: \emph{a crossing at \(k=X\)} (reported with its
calendar-day equivalent), \emph{already better at the smallest tested
\(k\)} (no crossing needed because the trained method starts ahead), and
\emph{no crossing within the tested range} (the trained method remains
behind throughout every tested \(k\), which is a distinct and, for the
zero-shot model's defense, more informative outcome than a numerical
breakpoint would be).
